Let \(M\in\mathfrak M_{\mathrm{SLCC}}\) and define
\[
T=\bigl(S,Z,(A_c)_{c\in\mathcal C},(Y_c)_{c\in\mathcal C}\bigr),\qquad
w_t=\Pr_M(T=t).
\]
The four unit-level restrictions imply the four membership conditions in \(\text{def:legal-response-types}\), so \(T\in\mathcal T_{\mathrm{SLCC}}\) almost surely.  Thus \(w\in\Delta_{\mathrm{SLCC}}\).

Fix \(o=(j,s,a,y,\ell)\in\mathcal O\).  Randomized assignment and the known arm law give \(\Pr(D=j,T=t)=\rho_jw_t\).  Conditional on \(D=j,T=t\), the common-response assumptions give \(A=a_{c_j(s)}(t)\) and \(Y=y_{c_j(s)}(t)\) on \(S=s\).  Revelation and masking give \(\ell=z_t\) when \(a=1\) and \(\ell=\star\) when \(a=0\).  Hence, by the incidence definition,
\[
\Pr_M(O_i=o)
=\sum_t\rho_jw_t\mathbf 1\!\left\{s_t=s,\ a_{c_j(s)}(t)=a,\ y_{c_j(s)}(t)=y,\
\ell=z_t\text{ if }a=1,\ \ell=\star\text{ if }a=0\right\}
=\sum_tB_{o,t}w_t.
\]
Therefore \(p=Bw\).  The target definition gives
\[
\tau_e(M)=\mathbb E_M[Y_{c_e(S)}]
=\sum_tw_ty_{c_e(s_t)}(t)=h_e^{\top}w.
\]

Conversely, fix \(w\in\Delta_{\mathrm{SLCC}}\).  Draw \(T=t\) with probability \(w_t\), independently draw \(D=j\) with probability \(\rho_j\), and define \(S,Z,(A_c),(Y_c)\) by the coordinates of \(T\).  Under policy \(j\), set \(A=A_{c_j(S)}\), \(Y=Y_{c_j(S)}\), \(R=A\), and set \(Z_{\mathrm{obs}}=Z\) on \(R=1\) and \(Z_{\mathrm{obs}}=\star\) on \(R=0\).  This construction satisfies randomized assignment, the known arm law, both common-response restrictions, revelation, and masking.  Every supported type is legal, so all four outcome restrictions hold.  Consequently \(M_w\in\mathfrak M_{\mathrm{SLCC}}\).  The preceding calculation gives observable law \(Bw\) and target \(h_e^{\top}w\), proving both inclusions.